% !PRINTED 95
\weberpage{95}

Wenn man zwei Colonnen in $A$ oder in $B$ vertauscht, so ändern sich die Elemente $c_i^{(h)}$ nicht, sondern vertauschen sich nur unter einander. Wenn man aber zwei Zeilen in $A$ oder in $B$ vertauscht, so ändern sich die $c_i^{(h)}$, indem die Factoren in den einzelnen Producten der Summe anders zusammengefasst werden; wenn man aber \gesperrt{entsprechende} Vertauschungen in den Zeilen von $A$ und von $B$ gleichzeitig vornimmt, so bleiben die $c_i^{(h)}$ ungeändert, weil dadurch nur die einzelnen Glieder der Summe vertauscht werden.

Indem man in $A$ oder in $B$ oder in beiden zugleich die Zeilen zu Colonnen macht, erhält man noch drei verschiedene Arten für die Bildung des Products zweier Determinanten in Determinantenform. Letzteres kann man auch so ausdrücken: \gesperrt{Um das Product zweier Determinanten zu bilden, kann man die Elemente der einzelnen Zeilen oder Colonnen des einen Factors mit den entsprechenden Elementen der Zeilen oder der Colonnen des anderen Factors multipliciren und die Producte addiren, und diese Productsummen als Elemente einer neuen Determinante auffassen.}

Auf ein Product zweier Determinanten mit verschiedener Elementenzahl lässt sich die Multiplicationsregel dadurch anwenden, dass man die Determinante mit geringerer Reihenzahl durch den Satz §. 22, IX. in eine andere mit mehr Reihen verwandelt.

\begin{center}
{\large \S.\ 28.}

\vspace{0.5em}
{\large\gesperrt{Determinanten der Unterdeterminanten.}}
\end{center}

Wir machen hier gleich eine Anwendung von dem Multiplicationsgesetz der Determinanten. Es sei wie bisher:
\[
A=\begin{vmatrix}
 a_1^{(1)}&a_2^{(1)}&\cdots&a_n^{(1)}\\
 a_1^{(2)}&a_2^{(2)}&\cdots&a_n^{(2)}\\
 \cdots&\cdots&&\cdots\\
 a_1^{(n)}&a_2^{(n)}&\cdots&a_n^{(n)}
\end{vmatrix},\tag{1}
\]
und
%CONT>>
